% Source: https://artofproblemsolving.com/community/c3h2380267_weekly_very_hard_problem
% Stable identity: aops-topic-2380267-q01
Since $x=0$ is not a root, divide by $x^2$:
\[x^2+x+1+\dfrac1x+\dfrac1{x^2}=0\iff\left(x+\dfrac1x\right)^2+\left(x+\dfrac1x\right)-1=0,\]
using $x^2+\dfrac1{x^2}=\left(x+\dfrac1x\right)^2-2$. So $t=x+\dfrac1x$ satisfies $t^2+t-1=0$, i.e. $t=\dfrac{-1\pm\sqrt5}{2}$. For each $t$, $x$ solves $x^2-tx+1=0$, so $x=\dfrac{t\pm\sqrt{t^2-4}}{2}$.
\begin{itemize}
\item $t=\dfrac{-1+\sqrt5}{2}$: $t^2=\dfrac{3-\sqrt5}{2}$ and $t^2-4=-\dfrac{5+\sqrt5}{2}$, so $x=\dfrac{-1+\sqrt5\pm i\sqrt{10+2\sqrt5}}{4}$.
\item $t=\dfrac{-1-\sqrt5}{2}$: $t^2=\dfrac{3+\sqrt5}{2}$ and $t^2-4=-\dfrac{5-\sqrt5}{2}$, so $x=\dfrac{-1-\sqrt5\pm i\sqrt{10-2\sqrt5}}{4}$.
\end{itemize}
(Here $\dfrac12\sqrt{\dfrac{5\pm\sqrt5}{2}}=\dfrac{\sqrt{10\pm2\sqrt5}}{4}$.) The four roots are
\[\boxed{x=\dfrac{-1+\sqrt5\pm i\sqrt{10+2\sqrt5}}{4},\qquad x=\dfrac{-1-\sqrt5\pm i\sqrt{10-2\sqrt5}}{4}}.\]

\emph{Another view.} Since $(x-1)\left(x^4+x^3+x^2+x+1\right)=x^5-1$, the roots are the fifth roots of unity other than $1$: $\cos72^\circ k+i\sin72^\circ k$ for $k=1,2,3,4$. Comparing real parts gives $\cos72^\circ=\dfrac{\sqrt5-1}{4}$ and $\cos144^\circ=-\dfrac{\sqrt5+1}{4}$.
